\chapter[FCC Condensate Constraint and Defect Impedance]{Interstitial-Condensate Constraint and Defect Impedance}
\label{app:condensate_impedance}

\vspace{1em}
\begin{center}
    \textbf{Abstract}
\end{center}
\begin{quote}
\noindent This appendix establishes the mechanical response of the interstitial condensate within the active FCC vacuum. It defines the incrementally incompressible coherent branch and separates the singular tangent volume constraint from the finite gravitational yield pressure. It then formulates the localized relative-flow action of a topological defect, constructs the mode-resolved added-inertia functional, and defines the microscopic impedance projection as a boundary-value functional with the dimensions of an inverse penetration length. The surface-to-volume relation is a reduced constitutive matching between the normal-coupled interface representation and the solved condensate response; it does not by itself establish a universal intrinsic condensate density. Canonical defect-flow solutions and the numerical evaluation of the projection remain open boundary-value problems.
\end{quote}
\vspace{2em}

\section{Scope and Constitutive Setting}

The active vacuum is a space-filling dual-phase medium composed of the solid FCC scaffold and the connected interstitial condensate. The primitive-cell volume fractions satisfy
\begin{equation}
    \beta_{\mathrm{geo}} = \frac{\pi}{3\sqrt{2}}, \qquad \beta_{\mathrm{IC}} = 1-\beta_{\mathrm{geo}}, \qquad \beta_{\mathrm{geo}} + \beta_{\mathrm{IC}} = 1.
\end{equation}

The FCC nodes occupy the close-packed solid fraction, while the interstitial condensate continuously fills the complementary octahedral and tetrahedral regions. No unoccupied third phase is introduced within the active FCC cell.

The analysis is restricted to the active phase,
\begin{equation}
    0 \leq \alpha_{\mathrm{G}} < 1, \qquad \dot{\alpha}_{\mathrm{G}} \approx 0,
\end{equation}
where $\alpha_{\mathrm{G}}$ specifies a gravitationally preloaded background that is quasistatic over the duration of the local disturbance. It is not itself treated as a propagating perturbation in the coherent branch.

The interstitial condensate has two distinct constitutive roles. In the coherent branch, it is carried with the FCC scaffold without relative flow and supplies the incremental pressure response that constrains local volumetric change. Under differential motion, it may redistribute relative to the scaffold and produce a conservative hydrodynamic reaction, together with any separately resolved dissipative impedance. These roles must not be combined into a single undifferentiated density or kinetic coefficient.

Let $\mathbf{u}(\mathbf{x},t)$ denote the homogenized FCC scaffold displacement and $\mathbf{w}(\mathbf{x},t)$ the interstitial-condensate displacement. Their relative velocity is
\begin{equation}
    \mathbf{v}_{\mathrm{rel}} = \dot{\mathbf{w}} - \dot{\mathbf{u}}.
\end{equation}

The coherent and localized branches are distinguished by
\begin{equation}
    \mathbf{v}_{\mathrm{rel}} = \mathbf{0} \quad \mathrm{(coherent)}, \qquad \mathbf{v}_{\mathrm{rel}} \neq \mathbf{0} \quad \mathrm{(localized\ relative\ flow)}.
\end{equation}

In the coherent branch, entrainment fixes the relative kinematics of the two phases but does not by itself establish an independently additive condensate contribution to the primitive-cell inertia. In the localized branch, nonzero relative motion generates the mode-dependent hydrodynamic impedance developed later.

Finite gravitational loading must be distinguished from the incremental constitutive response. The background may sustain finite pressure and approach the FCC saturation threshold as $\alpha_{\mathrm{G}}\rightarrow1^{-}$ while the coherent perturbation remains incrementally volume preserving. The resulting constraint does not prohibit quasistatic background compaction, nonlinear lattice rearrangement, or the FCC-to-SC phase transition.

The incremental actions are written about the current gravitationally preloaded configuration. Finite-deformation geometric-stiffness and acoustoelastic corrections are outside the retained tangent order. The localized construction later defines an equivalent mode-resolved impedance density whose numerical value and possible universality require an explicit defect-flow solution.

\section{Coherent Scaffold--Condensate Kinematics}
\label{sec:AP_coherent_kinematics}

In the entrained coherent branch, $\mathbf{w}=\mathbf{u}$ up to an irrelevant time-independent reference displacement. Consequently,
\begin{equation}
    \dot{\mathbf{w}} = \dot{\mathbf{u}}, \qquad \mathbf{v}_{\mathrm{rel}} = \mathbf{0}.
\end{equation}

The associated infinitesimal strain and volumetric strain are
\begin{equation}
    \varepsilon_{ij} = \frac{1}{2}\left(\partial_{i}u_{j}+\partial_{j}u_{i}\right), \qquad \varepsilon_{\mathrm{V}} = \mathrm{tr}\,\boldsymbol{\varepsilon} = \nabla\cdot\mathbf{u}.
\end{equation}

Let $\rho_{\mathrm{eff}}^{(0)}$ denote the unstrained inertia coefficient assigned to the selected coherent transfer coordinate. At the primitive-cell microscopic scale established in \hyperref[app:fcc_microphysics]{Appendix~\ref*{app:fcc_microphysics}},
\begin{equation}
    \rho_{\mathrm{eff}}^{(0)} = \rho_{\mathrm{sh}}^{(\mathrm{micro},0)} = \beta_{\mathrm{geo}}\rho_{\mathrm{D}}.
\end{equation}

This is the nodal scaffold mass averaged over the complete primitive-cell volume. Condensate entrainment fixes the coherent relative kinematics but does not independently add a condensate inertia to $\rho_{\mathrm{eff}}^{(0)}$. Any additional coherent condensate kinetic coefficient requires an explicit coupled two-field kinetic action.

The gravitational endpoint-transfer constraint gives the referred inertia
\begin{equation}
    \rho_{\mathrm{eff}}(\alpha_{\mathrm{G}}) = \frac{\rho_{\mathrm{eff}}^{(0)}}{\left(1-\alpha_{\mathrm{G}}\right)^4},
\end{equation}
and the coherent kinetic-energy density is
\begin{equation}
    \mathcal{T}_{\mathrm{coh}} = \frac{1}{2}\rho_{\mathrm{eff}}(\alpha_{\mathrm{G}})\dot{\mathbf{u}}\cdot\dot{\mathbf{u}}.
\end{equation}

The coefficient $\rho_{\mathrm{eff}}(\alpha_{\mathrm{G}})$ is a referred inertia of the selected coordinate, not a literal equilibrium density of the complete dual-phase vacuum. The localized condensate impedance density introduced later belongs to a distinct relative-flow coordinate and must not be added independently to this coherent kinetic term.

Let $K_{\mathrm{sh}}^{(0)}$ denote the shear coefficient paired with $\rho_{\mathrm{eff}}^{(0)}$ at the same descriptive scale:
\begin{equation}
    K_{\mathrm{sh}}^{(0)} = \rho_{\mathrm{eff}}^{(0)}c_{\mathrm{E}}^2.
\end{equation}

At the primitive-cell microscopic scale,
\begin{equation}
    K_{\mathrm{sh}}^{(0)} = K_{\mathrm{sh}}^{(\mathrm{micro},0)} = \rho_{\mathrm{sh}}^{(\mathrm{micro},0)}c_{\mathrm{E}}^2.
\end{equation}

At the retained tangent order,
\begin{equation}
    K_{\mathrm{sh}}(\alpha_{\mathrm{G}}) = K_{\mathrm{sh}}^{(0)}, \qquad 0 \leq \alpha_{\mathrm{G}} < 1.
\end{equation}

Before imposing the coherent volume constraint, introduce the finite diagnostic regularization coefficient $B_{\mathrm{IC}}$. The regularized elastic-energy density and isotropic tangent tensor are
\begin{equation}
    \mathcal{W}_{B} = K_{\mathrm{sh}}^{(0)}\varepsilon_{ij}\varepsilon_{ij} + \frac{1}{2}B_{\mathrm{IC}}\varepsilon_{\mathrm{V}}^2,
\end{equation}
and
\begin{equation}
    C_{ijkl}^{(B)} = B_{\mathrm{IC}}\delta_{ij}\delta_{kl} + K_{\mathrm{sh}}^{(0)}\left(\delta_{ik}\delta_{jl}+\delta_{il}\delta_{jk}\right).
\end{equation}

For propagation direction
\begin{equation}
    \hat{\mathbf{s}} = \frac{\mathbf{k}}{|\mathbf{k}|},
\end{equation}
the acoustic tensor is
\begin{equation}
    \Gamma_{ij} = C_{ikjl}^{(B)}s_{k}s_{l} = K_{\mathrm{sh}}^{(0)}\delta_{ij} + \left(B_{\mathrm{IC}}+K_{\mathrm{sh}}^{(0)}\right)s_{i}s_{j}.
\end{equation}

Its translational eigenvalues are
\begin{equation}
    \lambda_{\mathrm{T}} = K_{\mathrm{sh}}^{(0)}, \qquad \lambda_{\mathrm{L}} = B_{\mathrm{IC}} + 2K_{\mathrm{sh}}^{(0)}.
\end{equation}

The corresponding total bulk coefficient is
\begin{equation}
    K_{\mathrm{bulk}}^{(B)} = B_{\mathrm{IC}} + \frac{2}{3}K_{\mathrm{sh}}^{(0)}.
\end{equation}

The coefficient $B_{\mathrm{IC}}$ is introduced only to expose the hypothetical finite longitudinal acoustic branch. It is not identified with the finite FCC saturation-pressure scale; the tangent volume constraint and the nonlinear phase-transition pressure are distinct constitutive quantities.

\section{Finite-Compressibility Regularization}
\label{sec:AP_finite_regularization}

The quadratic action of the regularized coherent branch is
\begin{equation}
    S_{B}^{(2)} = \int \left[ \frac{1}{2}\rho_{\mathrm{eff}}(\alpha_{\mathrm{G}})\dot{u}_{i}\dot{u}_{i} - K_{\mathrm{sh}}^{(0)}\varepsilon_{ij}\varepsilon_{ij} - \frac{1}{2}B_{\mathrm{IC}}\varepsilon_{\mathrm{V}}^2 \right] \mathrm{d}^3x\,\mathrm{d}t.
\end{equation}

At the presently derived primitive-cell level, the kinetic coefficient reduces to the FCC scaffold inertia normalization. The condensate enters the coherent reduction through entrainment and the volumetric pressure response rather than through an independently calibrated additive kinetic density.

Variation with respect to $\mathbf{u}$ gives
\begin{equation}
    \sigma_{ij} = 2K_{\mathrm{sh}}^{(0)}\varepsilon_{ij} + B_{\mathrm{IC}}\varepsilon_{\mathrm{V}}\delta_{ij},
\end{equation}
and
\begin{equation}
    \rho_{\mathrm{eff}}(\alpha_{\mathrm{G}})\ddot{\mathbf{u}} = K_{\mathrm{sh}}^{(0)}\nabla^2\mathbf{u} + \left(B_{\mathrm{IC}}+K_{\mathrm{sh}}^{(0)}\right)\nabla\left(\nabla\cdot\mathbf{u}\right).
\end{equation}

For the plane-wave field
\begin{equation}
    \mathbf{u} = \mathbf{e}\exp\left[i\left(\mathbf{k}\cdot\mathbf{x}-\omega t\right)\right],
\end{equation}
the transverse and longitudinal branches satisfy
\begin{equation}
    \omega_{\mathrm{T}}^2 = \frac{K_{\mathrm{sh}}^{(0)}}{\rho_{\mathrm{eff}}(\alpha_{\mathrm{G}})}k^2, \qquad \omega_{\mathrm{L}}^2 = \frac{B_{\mathrm{IC}}+2K_{\mathrm{sh}}^{(0)}}{\rho_{\mathrm{eff}}(\alpha_{\mathrm{G}})}k^2.
\end{equation}

In the unstrained active vacuum, $\alpha_{\mathrm{G}}=0$, so
\begin{equation}
    c_{\mathrm{T}}(0) = c_{\mathrm{E}}, \qquad c_{\mathrm{L}}^2(0) = c_{\mathrm{E}}^2\left(\frac{B_{\mathrm{IC}}}{K_{\mathrm{sh}}^{(0)}}+2\right).
\end{equation}

Mechanical stability requires
\begin{equation}
    K_{\mathrm{bulk}}^{(B)} > 0 \qquad \Longrightarrow \qquad \frac{B_{\mathrm{IC}}}{K_{\mathrm{sh}}^{(0)}} > -\frac{2}{3},
\end{equation}
and therefore
\begin{equation}
    c_{\mathrm{L}}^2(0) > \frac{4}{3}c_{\mathrm{E}}^2, \qquad c_{\mathrm{L}}(0) > \sqrt{\frac{4}{3}}\,c_{\mathrm{E}}.
\end{equation}

For $B_{\mathrm{IC}} \geq 0$,
\begin{equation}
    c_{\mathrm{L}}(0) \geq \sqrt{2}\,c_{\mathrm{E}}.
\end{equation}

Thus an ordinary finite isotropic elastic model cannot simultaneously satisfy $c_{\mathrm{T}}=c_{\mathrm{E}}$, retain $c_{\mathrm{E}}$ as the active-phase propagation ceiling, and support an independent stable finite-frequency longitudinal acoustic branch generated by $B_{\mathrm{IC}}$.

The finite-$B_{\mathrm{IC}}$ theory therefore serves only as a diagnostic regularization. It does not uniquely derive incompressibility from stability; rather, it motivates a reduced coherent constitutive branch in which local volumetric fluctuations are excluded from the finite-energy acoustic spectrum.

\section{Incrementally Incompressible Coherent Branch}
\label{sec:AP_incompressible_branch}

The coherent active-FCC branch is defined by the singular constitutive reduction
\begin{equation}
    B_{\mathrm{IC}} \rightarrow \infty.
\end{equation}

Finite perturbation energy then requires
\begin{equation}
    \frac{1}{2}B_{\mathrm{IC}}\varepsilon_{\mathrm{V}}^2 < \infty \qquad \Longrightarrow \qquad \varepsilon_{\mathrm{V}} \rightarrow 0,
\end{equation}
so that
\begin{equation}
    \nabla\cdot\mathbf{u} = 0.
\end{equation}

This is an incremental constraint about a prescribed gravitationally loaded background. It neither requires zero finite background compaction nor prohibits approach to the nonlinear saturation boundary.

Introduce the incremental condensate pressure $\delta p_{\mathrm{IC}}(\mathbf{x},t)$ as the Lagrange multiplier enforcing the volume constraint. The reduced coherent action is
\begin{equation}
    S_{\mathrm{coh}}^{(2)} = \int \left[ \frac{1}{2}\rho_{\mathrm{eff}}(\alpha_{\mathrm{G}})\dot{u}_{i}\dot{u}_{i} - K_{\mathrm{sh}}^{(0)}\varepsilon_{ij}\varepsilon_{ij} + \delta p_{\mathrm{IC}}\nabla\cdot\mathbf{u} \right] \mathrm{d}^3x\,\mathrm{d}t.
\end{equation}

At the primitive-cell microscopic scale,
\begin{equation}
    \rho_{\mathrm{eff}}^{(0)} = \rho_{\mathrm{sh}}^{(\mathrm{micro},0)}.
\end{equation}
The condensate enters this reduced action through entrainment and the pressure constraint; no independently additive coherent condensate kinetic density is introduced at the present constitutive order.

A constituent-continuity check is consistent with the same constraint. Let $\varrho_{\mathrm{IC}}$ denote a formal local condensate constituent-density field and define
\begin{equation}
    \widetilde{\rho}_{\mathrm{IC}} = \beta_{\mathrm{IC}}\varrho_{\mathrm{IC}}.
\end{equation}

The quantity $\varrho_{\mathrm{IC}}$ is used only for this continuity construction. No independently calibrated equilibrium value or additive coherent inertia is inferred from it.

The local conservation law is
\begin{equation}
    \frac{\partial\widetilde{\rho}_{\mathrm{IC}}}{\partial t} + \nabla\cdot\left(\widetilde{\rho}_{\mathrm{IC}}\mathbf{v}_{\mathrm{IC}}\right) = 0.
\end{equation}

For the entrained branch,
\begin{equation}
    \mathbf{v}_{\mathrm{IC}} = \dot{\mathbf{u}},
\end{equation}
and, over the incremental disturbance,
\begin{equation}
    \frac{D\beta_{\mathrm{IC}}}{Dt} = 0, \qquad \frac{D\varrho_{\mathrm{IC}}}{Dt} = 0.
\end{equation}

The continuity equation then gives
\begin{equation}
    \nabla\cdot\dot{\mathbf{u}} = 0,
\end{equation}
which integrates to $\nabla\cdot\mathbf{u}=0$ for an initially volume-preserving incremental reference state. This is a consistency check on the selected coherent branch, not an independent derivation of incompressibility. The fields $\varrho_{\mathrm{IC}}$ and $\widetilde{\rho}_{\mathrm{IC}}$ remain distinct from the later defect impedance density $\rho_{\mathrm{IC}}^{(\mathrm{imp})}$.

Variation of the constrained action gives
\begin{equation}
    \delta\sigma_{ij} = -\delta p_{\mathrm{IC}}\delta_{ij} + 2K_{\mathrm{sh}}^{(0)}\varepsilon_{ij},
\end{equation}
and
\begin{equation}
    \rho_{\mathrm{eff}}(\alpha_{\mathrm{G}})\ddot{\mathbf{u}} = K_{\mathrm{sh}}^{(0)}\nabla^2\mathbf{u} - \nabla\delta p_{\mathrm{IC}}, \qquad \nabla\cdot\mathbf{u} = 0.
\end{equation}

In a source-free homogeneous region,
\begin{equation}
    \nabla^2\delta p_{\mathrm{IC}} = 0.
\end{equation}

With an applied incremental force density $\mathbf{f}$,
\begin{equation}
    \rho_{\mathrm{eff}}(\alpha_{\mathrm{G}})\ddot{\mathbf{u}} = K_{\mathrm{sh}}^{(0)}\nabla^2\mathbf{u} - \nabla\delta p_{\mathrm{IC}} + \mathbf{f},
\end{equation}
and
\begin{equation}
    \nabla^2\delta p_{\mathrm{IC}} = \nabla\cdot\mathbf{f}.
\end{equation}

The elliptic pressure equation is the compatibility condition of the constrained branch. It does not define an independently propagating field or an additional channel of instantaneous information.

For a plane wave, incompressibility requires
\begin{equation}
    \mathbf{k}\cdot\mathbf{e} = 0.
\end{equation}

The longitudinal mode of the finite-$B_{\mathrm{IC}}$ regularization satisfies
\begin{equation}
    \omega_{\mathrm{L}}^2 = \frac{B_{\mathrm{IC}}+2K_{\mathrm{sh}}^{(0)}}{\rho_{\mathrm{eff}}(\alpha_{\mathrm{G}})}k^2,
\end{equation}
so for fixed $k \neq 0$,
\begin{equation}
    B_{\mathrm{IC}} \rightarrow \infty \qquad \Longrightarrow \qquad \omega_{\mathrm{L}} \rightarrow \infty.
\end{equation}

The longitudinal translational excitation is therefore removed from the finite-frequency coherent spectrum rather than converted into a superluminal propagating signal.

For either admissible transverse polarization,
\begin{equation}
    \omega_{\mathrm{T}}^2 = \frac{K_{\mathrm{sh}}^{(0)}}{\rho_{\mathrm{eff}}(\alpha_{\mathrm{G}})}k^2 + \mathcal{O}\left(k^4d_{\mathrm{nn}}^2\right).
\end{equation}

Using
\begin{equation}
    K_{\mathrm{sh}}^{(0)} = \rho_{\mathrm{eff}}^{(0)}c_{\mathrm{E}}^2, \qquad \rho_{\mathrm{eff}}(\alpha_{\mathrm{G}}) = \frac{\rho_{\mathrm{eff}}^{(0)}}{\left(1-\alpha_{\mathrm{G}}\right)^4},
\end{equation}
one obtains
\begin{equation}
    \omega_{\mathrm{T}}^2 = c_{\mathrm{E}}^2\left(1-\alpha_{\mathrm{G}}\right)^4 k^2 + \mathcal{O}\left(k^4d_{\mathrm{nn}}^2\right),
\end{equation}
and hence
\begin{equation}
    c_{\mathrm{P}} = c_{\mathrm{E}}\left(1-\alpha_{\mathrm{G}}\right)^2.
\end{equation}

The translational acoustic sector of the coherent active vacuum therefore contains two finite-frequency transverse polarizations. This statement does not determine the independent dynamical status of internal phase, microrotational, torsional, chiral, or localized-defect modes.

\section{Static Longitudinal Pressure and the Finite FCC Yield Scale}
\label{sec:AP_static_pressure_yield}

Incremental incompressibility does not imply the absence of hydrostatic stress. Decompose the interstitial-condensate pressure as
\begin{equation}
    p_{\mathrm{IC}}(\mathbf{x},t) = \overline{p}_{\mathrm{IC}}(\mathbf{x}) + \delta p_{\mathrm{IC}}(\mathbf{x},t),
\end{equation}
where $\overline{p}_{\mathrm{IC}}$ supports the quasistatic gravitational preload and $\delta p_{\mathrm{IC}}$ enforces the incremental coherent volume constraint.

The local gravitational loading is
\begin{equation}
    \tau_{\mathrm{L}} = \alpha_{\mathrm{G}}\tau_{\mathrm{D}},
\end{equation}
where $\tau_{\mathrm{D}}$ is the active Soft Floor surface-tension scale. The corresponding finite volumetric yield-pressure scale is
\begin{equation}
    \tau_{\mathrm{D,bulk}} \equiv \frac{\tau_{\mathrm{D}}}{r_{\mathrm{D}}}.
\end{equation}

The background condensate pressure is therefore
\begin{equation}
    \overline{p}_{\mathrm{IC}} = \frac{\tau_{\mathrm{L}}}{r_{\mathrm{D}}} = \alpha_{\mathrm{G}}\tau_{\mathrm{D,bulk}},
\end{equation}
or equivalently
\begin{equation}
    \alpha_{\mathrm{G}} = \frac{\overline{p}_{\mathrm{IC}}}{\tau_{\mathrm{D,bulk}}}.
\end{equation}

Within the active phase,
\begin{equation}
    0 \leq \overline{p}_{\mathrm{IC}} < \tau_{\mathrm{D,bulk}}.
\end{equation}
At $\alpha_{\mathrm{G}}=0$, $\overline{p}_{\mathrm{IC}}=0$, while
\begin{equation}
    \alpha_{\mathrm{G}} \rightarrow 1^{-} \qquad \Longrightarrow \qquad \overline{p}_{\mathrm{IC}} \rightarrow \tau_{\mathrm{D,bulk}}.
\end{equation}

At equality, the active FCC constitutive coordinates terminate and the local structure enters the saturated phase-transition regime.

The tangent incompressibility parameter and the finite saturation pressure have distinct roles:
\begin{equation}
    B_{\mathrm{IC}} \rightarrow \infty, \qquad \nabla\cdot\mathbf{u} = 0,
\end{equation}
whereas
\begin{equation}
    \tau_{\mathrm{D,bulk}} < \infty
\end{equation}
sets the maximum quasistatic pressure sustainable by the active FCC phase. Therefore,
\begin{equation}
    B_{\mathrm{IC}} \neq \tau_{\mathrm{D,bulk}}.
\end{equation}

The total incremental stress about the preloaded state is
\begin{equation}
    \delta\sigma_{ij} = -\delta p_{\mathrm{IC}}\delta_{ij} + 2K_{\mathrm{sh}}^{(0)}\varepsilon_{ij},
\end{equation}
and the corresponding total stress may be written
\begin{equation}
    \sigma_{ij} = -\overline{p}_{\mathrm{IC}}\delta_{ij} - \delta p_{\mathrm{IC}}\delta_{ij} + 2K_{\mathrm{sh}}^{(0)}\varepsilon_{ij}.
\end{equation}

The background term satisfies the quasistatic equilibrium equations independently of the incremental disturbance. Finite-deformation geometric-stiffness corrections from the preload are outside the retained tangent order.

Using
\begin{equation}
    \alpha_{\mathrm{G}} = \frac{\overline{p}_{\mathrm{IC}}}{\tau_{\mathrm{D,bulk}}},
\end{equation}
the locally uniform endpoint-transfer factor becomes
\begin{equation}
    \mathcal{C} = \left(1-\frac{\overline{p}_{\mathrm{IC}}}{\tau_{\mathrm{D,bulk}}}\right)^2,
\end{equation}
and the referred coherent inertia is
\begin{equation}
    \rho_{\mathrm{eff}} = \frac{\rho_{\mathrm{eff}}^{(0)}}{\left(1-\overline{p}_{\mathrm{IC}}/\tau_{\mathrm{D,bulk}}\right)^4}.
\end{equation}

The transverse propagation speed is therefore
\begin{equation}
    c_{\mathrm{P}} = c_{\mathrm{E}}\left(1-\frac{\overline{p}_{\mathrm{IC}}}{\tau_{\mathrm{D,bulk}}}\right)^2.
\end{equation}

The static longitudinal pressure and the dynamic transverse propagation law are thus coupled through the occupied active capacity while remaining mechanically distinct.

\section{Phase-Sector Status and Translational Acoustic Mode Count}
\label{sec:AP_phase_status}

The objective contact action developed in \hyperref[app:fcc_microphysics]{Appendix~\ref*{app:fcc_microphysics}} contains the internal phase difference
\begin{equation}
    \Delta\varphi_{ij} = \varphi_{j}-\varphi_{i},
\end{equation}
with local phase energy
\begin{equation}
    \mathcal{U}_{\varphi}^{(ij)} = \kappa_{\varphi}r_{\mathrm{D}}^2\left[1-\cos(\Delta\varphi_{ij})\right].
\end{equation}

In the harmonic regime,
\begin{equation}
    \mathcal{U}_{\varphi}^{(ij)} = \frac{1}{2}\kappa_{\varphi}r_{\mathrm{D}}^2(\Delta\varphi_{ij})^2 + \mathcal{O}\left((\Delta\varphi_{ij})^4\right).
\end{equation}

Because the contact energy depends only on phase differences, the uniform transformation
\begin{equation}
    \varphi_{i} \rightarrow \varphi_{i}+\varphi_{0}
\end{equation}
is a null direction of the phase-difference potential.

This fact does not establish an independent propagating bulk phase mode. Such a mode requires a separately derived phase kinetic sector, which is not supplied by the present constitutive reduction. The phase coordinate may therefore represent an independent propagating field, a constrained internal coordinate, or a defect-confined degree of freedom depending on the completed microscopic dynamics. That classification remains open.

The mode count established here applies only to the translational acoustic sector. The coherent constraint
\begin{equation}
    \nabla\cdot\mathbf{u} = 0
\end{equation}
removes the longitudinal translational polarization and leaves two finite-frequency transverse translational polarizations.

For either transverse polarization,
\begin{equation}
    \mathbf{k}\cdot\mathbf{e}_{\mathrm{T}} = 0,
\end{equation}
with
\begin{equation}
    \omega_{\mathrm{T}}^2 = c_{\mathrm{E}}^2(1-\alpha_{\mathrm{G}})^4k^2 + \mathcal{O}\left(k^4d_{\mathrm{nn}}^2\right).
\end{equation}

Under the leading isotropy condition
\begin{equation}
    \kappa_{\mathrm{n}} = \kappa_{\mathrm{t}},
\end{equation}
the two transverse translational modes are degenerate at leading long-wavelength order. Finite-wave-number cubic anisotropy, chiral structure, nonuniform loading, higher-gradient contact terms, or additional phase--rotation coupling may lift that degeneracy.

Microrotational, torsional, internal-phase, chiral, and localized-defect modes remain outside this translational acoustic mode count.

\section{Localized Defect--Condensate Relative-Flow Kinematics}
\label{sec:AP_localized_kinematics}

The coherent branch satisfies
\begin{equation}
    \mathbf{v}_{\mathrm{rel}} = \dot{\mathbf{w}}-\dot{\mathbf{u}} = \mathbf{0}.
\end{equation}

A localized topological defect defines a distinct kinematic sector in which translation, boundary deformation, or internal reconfiguration can induce condensate motion relative to the local FCC scaffold:
\begin{equation}
    \mathbf{v}_{\mathrm{rel}} \neq \mathbf{0}.
\end{equation}

Let $\Omega_{\mathrm{IC}}^{(\mathrm{def})}$ denote the connected interstitial-condensate domain influenced by the defect and $\partial\Omega_{\mathrm{def}}$ the effective defect--condensate interface.

A single localized mode is described by a generalized coordinate $q(t)$ with
\begin{equation}
    [q] = \mathrm{m}.
\end{equation}

This length-valued normalization ensures that the boundary mode is dimensionless. Dimensionless internal coordinates require a separate coordinate metric and are not included in this scalar translational normalization.

Represent a material point on the defect boundary by
\begin{equation}
    \mathbf{X}_{\partial} = \mathbf{X}_{\partial}(\boldsymbol{s},q),
\end{equation}
and define
\begin{equation}
    \mathbf{b}_{\partial} \equiv \frac{\partial\mathbf{X}_{\partial}}{\partial q}, \qquad [\mathbf{b}_{\partial}] = 1.
\end{equation}

The boundary velocity relative to the scaffold is
\begin{equation}
    \dot{\mathbf{X}}_{\partial}^{(\mathrm{rel})} = \dot{q}\mathbf{b}_{\partial}.
\end{equation}

For a linearized localized response,
\begin{equation}
    \mathbf{v}_{\mathrm{rel}}(\mathbf{x},t) = \dot{q}(t)\boldsymbol{\Upsilon}(\mathbf{x}),
\end{equation}
where $\boldsymbol{\Upsilon}$ is the dimensionless condensate response to unit generalized defect velocity.

The field $\boldsymbol{\Upsilon}$ depends on defect geometry, interstitial connectivity, outer boundary conditions, and the selected mode. At conservative inviscid order, the no-penetration condition is
\begin{equation}
    \boldsymbol{\Upsilon}\cdot\hat{\mathbf{n}}_{\mathrm{d}} = \mathbf{b}_{\partial}\cdot\hat{\mathbf{n}}_{\mathrm{d}} \qquad \mathrm{on}\ \partial\Omega_{\mathrm{def}},
\end{equation}
where $\hat{\mathbf{n}}_{\mathrm{d}}$ points from the defect into the condensate.

Only the normal boundary velocity is transmitted by this baseline inviscid condition. Tangential defect motion requires an additional tangential traction, interfacial locking, or no-slip law to contribute to the condensate response.

The kinetic contribution generated by $\boldsymbol{\Upsilon}$ is therefore an added inertia of the localized defect coordinate, not an independent contribution to the coherent acoustic inertia.

\section[Defect-Flow Boundary-Value Problem]{Conservative Defect-Flow Boundary-Value Problem}
\label{sec:AP_defect_boundary_value_problem}

The localized interstitial-condensate response is incrementally incompressible:
\begin{equation}
    \nabla\cdot\boldsymbol{\Upsilon} = 0
\end{equation}
throughout $\Omega_{\mathrm{IC}}^{(\mathrm{def})}$.

On the minimum-energy zero-circulation branch,
\begin{equation}
    \nabla\times\boldsymbol{\Upsilon} = \mathbf{0},
\end{equation}
so that
\begin{equation}
    \boldsymbol{\Upsilon} = \nabla\Phi_{\mathrm{rel}}.
\end{equation}

Since $[\boldsymbol{\Upsilon}]=1$,
\begin{equation}
    [\Phi_{\mathrm{rel}}] = \mathrm{m}.
\end{equation}
The physical velocity potential is $\dot{q}\Phi_{\mathrm{rel}}$.

The reduced potential satisfies
\begin{equation}
    \nabla^2\Phi_{\mathrm{rel}} = 0
\end{equation}
within $\Omega_{\mathrm{IC}}^{(\mathrm{def})}$, with the inner boundary condition
\begin{equation}
    \frac{\partial\Phi_{\mathrm{rel}}}{\partial n_{\mathrm{d}}} = \mathbf{b}_{\partial}\cdot\hat{\mathbf{n}}_{\mathrm{d}} \qquad \mathrm{on}\ \partial\Omega_{\mathrm{def}}.
\end{equation}

For an isolated defect in an unbounded condensate domain,
\begin{equation}
    \nabla\Phi_{\mathrm{rel}} \rightarrow \mathbf{0} \qquad \mathrm{as}\ |\mathbf{x}| \rightarrow \infty.
\end{equation}

For a periodic FCC calculation, the corresponding periodic boundary conditions and cell-scale flux constraint replace the far-field condition.

Let $\partial\Omega_{\mathrm{out}}$ denote the outer boundary of the condensate domain. Flux compatibility requires
\begin{equation}
    -\int_{\partial\Omega_{\mathrm{def}}}\mathbf{b}_{\partial}\cdot\hat{\mathbf{n}}_{\mathrm{d}}\,\mathrm{d}A + \int_{\partial\Omega_{\mathrm{out}}}\boldsymbol{\Upsilon}\cdot\hat{\mathbf{n}}_{\mathrm{out}}\,\mathrm{d}A = 0.
\end{equation}

For an isolated defect with vanishing far-field flux,
\begin{equation}
    \int_{\partial\Omega_{\mathrm{def}}}\mathbf{b}_{\partial}\cdot\hat{\mathbf{n}}_{\mathrm{d}}\,\mathrm{d}A = 0.
\end{equation}

Rigid translation satisfies this condition because
\begin{equation}
    \int_{\partial\Omega_{\mathrm{def}}}\hat{\mathbf{n}}_{\mathrm{d}}\,\mathrm{d}A = \mathbf{0}.
\end{equation}

A purely radial breathing mode that changes the enclosed defect volume does not satisfy the isolated-domain condition unless its flux is balanced by outer-boundary flow, compensating scaffold deformation, or a coupled phase-conversion process. It is therefore not an admissible isolated perturbation of the strictly incompressible conservative branch.

The conservative response may equivalently be characterized as
\begin{equation}
    \boldsymbol{\Upsilon} = \operatorname*{arg\,min}_{\mathbf{v}} \int_{\Omega_{\mathrm{IC}}^{(\mathrm{def})}}|\mathbf{v}|^2\,\mathrm{d}V,
\end{equation}
subject to
\begin{equation}
    \nabla\cdot\mathbf{v} = 0, \qquad \mathbf{v}\cdot\hat{\mathbf{n}}_{\mathrm{d}} = \mathbf{b}_{\partial}\cdot\hat{\mathbf{n}}_{\mathrm{d}} \qquad \mathrm{on}\ \partial\Omega_{\mathrm{def}},
\end{equation}
together with the selected outer-boundary condition.

A circulation-carrying or chiral defect requires
\begin{equation}
    \boldsymbol{\Upsilon} = \nabla\Phi_{\mathrm{rel}} + \boldsymbol{\Upsilon}_{\mathrm{circ}},
\end{equation}
with
\begin{equation}
    \nabla\cdot\boldsymbol{\Upsilon}_{\mathrm{circ}} = 0, \qquad \nabla\times\boldsymbol{\Upsilon}_{\mathrm{circ}} \neq \mathbf{0}.
\end{equation}

Such branches lie outside the baseline zero-circulation impedance projection. For a prescribed admissible geometry, boundary mode, and outer condition, the Neumann problem determines $\Phi_{\mathrm{rel}}$ up to an additive constant on a simply connected domain. Multiply connected domains may additionally contain harmonic circulation sectors that require separate specification.

\section[Normal Boundary Impedance Projection]{Normal Boundary Weighting and Impedance Projection}
\label{sec:AP_normal_boundary_impedance}

The conservative boundary-value problem transmits the normal component of the defect-boundary velocity into the interstitial condensate. Define
\begin{equation}
    b_{\mathrm{n}}(\boldsymbol{s}) \equiv \mathbf{b}_{\partial}(\boldsymbol{s})\cdot\hat{\mathbf{n}}_{\mathrm{d}}(\boldsymbol{s}).
\end{equation}

The mode-weighted normal interface area is
\begin{equation}
    A_{\mathrm{imp}}^{(\mathrm{n})} \equiv \int_{\partial\Omega_{\mathrm{def}}}|b_{\mathrm{n}}|^2\,\mathrm{d}A, \qquad [A_{\mathrm{imp}}^{(\mathrm{n})}] = \mathrm{m}^2.
\end{equation}

The quantity $A_{\mathrm{imp}}^{(\mathrm{n})}$ is not generally the total geometric defect area. It is the normal-coupled modal area transmitted by the no-penetration condition; boundary regions with vanishing normal motion do not contribute.

The condensate participation volume is
\begin{equation}
    V_{\mathrm{imp}} \equiv \int_{\Omega_{\mathrm{IC}}^{(\mathrm{def})}}|\boldsymbol{\Upsilon}|^2\,\mathrm{d}V, \qquad [V_{\mathrm{imp}}] = \mathrm{m}^3.
\end{equation}

It is an inertia-weighted participation volume rather than the geometric volume of the entire connected interstitial domain.

Define the localized impedance projection by
\begin{equation}
    \label{eq:AP_defect_projection}
    \iota_{\mathrm{V}}^{(\mathrm{def})} \equiv \frac{A_{\mathrm{imp}}^{(\mathrm{n})}}{V_{\mathrm{imp}}}, \qquad [\iota_{\mathrm{V}}^{(\mathrm{def})}] = \mathrm{m}^{-1}.
\end{equation}

This quantity must be distinguished from the periodic FCC scaffold modal projection established in \hyperref[app:fcc_microphysics]{Appendix~\ref*{app:fcc_microphysics}},
\begin{equation}
    \iota_{\mathrm{tr}}^{(\mathrm{modal})} = \frac{A_{\mathrm{ct}}^{(\mathrm{modal})}}{V_{\mathrm{p}}} = \frac{\rho_{\mathrm{sh}}^{(\mathrm{micro},0)}}{\mu_{\mathrm{D}}}.
\end{equation}

The two projections satisfy
\begin{equation}
    \label{eq:AP_projection_distinction}
    \iota_{\mathrm{V}}^{(\mathrm{def})} \not\equiv \iota_{\mathrm{tr}}^{(\mathrm{modal})}.
\end{equation}

The quantity $\iota_{\mathrm{tr}}^{(\mathrm{modal})}$ belongs to the periodic scaffold normalization, whereas $\iota_{\mathrm{V}}^{(\mathrm{def})}$ belongs to a localized relative-condensate-flow boundary-value problem. Equality requires an independent demonstration that the corresponding mode functionals coincide.

In terms of the solved fields,
\begin{equation}
    \iota_{\mathrm{V}}^{(\mathrm{def})} = \frac{\displaystyle\int_{\partial\Omega_{\mathrm{def}}}|\mathbf{b}_{\partial}\cdot\hat{\mathbf{n}}_{\mathrm{d}}|^2\,\mathrm{d}A} {\displaystyle\int_{\Omega_{\mathrm{IC}}^{(\mathrm{def})}}|\boldsymbol{\Upsilon}|^2\,\mathrm{d}V}.
\end{equation}

The projection is invariant under a dimensionless rescaling $q'=a q$, since
\begin{equation}
    \mathbf{b}_{\partial}' = a^{-1}\mathbf{b}_{\partial}, \qquad \boldsymbol{\Upsilon}' = a^{-1}\boldsymbol{\Upsilon},
\end{equation}
implies
\begin{equation}
    A_{\mathrm{imp}}^{(\mathrm{n})\prime} = a^{-2}A_{\mathrm{imp}}^{(\mathrm{n})}, \qquad V_{\mathrm{imp}}' = a^{-2}V_{\mathrm{imp}}, \qquad \iota_{\mathrm{V}}^{(\mathrm{def})\prime} = \iota_{\mathrm{V}}^{(\mathrm{def})}.
\end{equation}

Define the effective impedance penetration length by
\begin{equation}
    \ell_{\mathrm{imp}}^{(\mathrm{def})} \equiv \frac{V_{\mathrm{imp}}}{A_{\mathrm{imp}}^{(\mathrm{n})}},
\end{equation}
so that
\begin{equation}
    \iota_{\mathrm{V}}^{(\mathrm{def})} = \frac{1}{\ell_{\mathrm{imp}}^{(\mathrm{def})}}.
\end{equation}

The quantity $\ell_{\mathrm{imp}}^{(\mathrm{def})}$ is a mode measure and need not equal the defect radius, shell thickness, or lattice spacing.

The empirical convention
\begin{equation}
    \iota_{\mathrm{V}} = 1\ \mathrm{m}^{-1}
\end{equation}
corresponds to
\begin{equation}
    \ell_{\mathrm{imp}} = 1\ \mathrm{m},
\end{equation}
and remains a calibration assumption until a specified defect-flow boundary-value problem produces the same mode-integral ratio.

\section{Mode-Resolved Condensate Impedance Density}
\label{sec:AP_mode_resolved_density}

Adopt the reduced normal-coupled surface normalization
\begin{equation}
    T_{\partial}^{(\mathrm{n})} \equiv \frac{1}{2}\mu_{\mathrm{D}}A_{\mathrm{imp}}^{(\mathrm{n})}\dot{q}^2.
\end{equation}

This is a constitutive interface representation of the condensate-coupled normal mode. The coefficient $\mu_{\mathrm{D}}$ serves as the baseline surface-inertia normalization of the coupled interface response; it is not identified with an intrinsic volumetric density of the interstitial condensate.

Represent the same reduced kinetic coefficient in volumetric form as
\begin{equation}
    T_{\mathrm{vol}} \equiv \frac{1}{2}\rho_{\mathrm{IC}}^{(\mathrm{imp})}V_{\mathrm{imp}}\dot{q}^2.
\end{equation}

Equivalence of the two representations gives
\begin{equation}
    \mu_{\mathrm{D}}A_{\mathrm{imp}}^{(\mathrm{n})} = \rho_{\mathrm{IC}}^{(\mathrm{imp})}V_{\mathrm{imp}},
\end{equation}
and therefore
\begin{equation}
    \rho_{\mathrm{IC}}^{(\mathrm{imp})} \equiv \mu_{\mathrm{D}}\iota_{\mathrm{V}}^{(\mathrm{def})} = \frac{\mu_{\mathrm{D}}}{\ell_{\mathrm{imp}}^{(\mathrm{def})}}.
\end{equation}

Its dimensions are
\begin{equation}
    [\rho_{\mathrm{IC}}^{(\mathrm{imp})}] = \mathrm{kg\,m^{-3}}.
\end{equation}

The quantity $\rho_{\mathrm{IC}}^{(\mathrm{imp})}$ is an equivalent mode-resolved impedance density associated with the localized relative-flow coordinate. It is neither automatically the intrinsic equilibrium density of the condensate nor an independently additive density of the coherent scaffold coordinate.

The corresponding conservative added-inertia coefficient is
\begin{equation}
    m_{\mathrm{add}} = \mu_{\mathrm{D}}A_{\mathrm{imp}}^{(\mathrm{n})} = \rho_{\mathrm{IC}}^{(\mathrm{imp})}V_{\mathrm{imp}}.
\end{equation}

This coefficient belongs exclusively to the localized relative-condensate-flow problem. It must not be added to the primitive-cell scaffold inertia $\rho_{\mathrm{sh}}^{(\mathrm{micro},0)}$ merely because the condensate is entrained in the coherent branch. Any independently additive coherent condensate kinetic coefficient requires an explicit coupled two-field kinetic action.

The localized projection may depend on geometry, topology, boundary mode, and outer condition,
\begin{equation}
    \iota_{\mathrm{V}}^{(\mathrm{def})} = \iota_{\mathrm{V}}^{(\mathrm{def})} \left( \mathrm{geometry}, \mathrm{topology}, \mathrm{boundary\ mode}, \mathrm{outer\ condition} \right).
\end{equation}

In a conservative frequency-dependent calculation, the response may define a real nonnegative projection $\iota_{\mathrm{V}}^{(\mathrm{def})}(\omega)$ with quasistatic limit
\begin{equation}
    \iota_{\mathrm{V},0}^{(\mathrm{def})} \equiv \lim_{\omega\rightarrow0}\iota_{\mathrm{V}}^{(\mathrm{def})}(\omega),
\end{equation}
provided the limit exists.

A dissipative complex impedance is a distinct response object and must be defined through a force--velocity relation such as
\begin{equation}
    \widehat{Q}_{\mathrm{def}}(\omega) = \mathcal{Z}_{\mathrm{def}}(\omega)\widehat{\dot{q}}(\omega).
\end{equation}

It cannot be obtained by assigning an imaginary part to the positive norm ratio $A_{\mathrm{imp}}^{(\mathrm{n})}/V_{\mathrm{imp}}$. Dissipative and radiative corrections therefore lie outside the conservative action retained here.

A universal condensate impedance density may be introduced only if $\iota_{\mathrm{V},0}^{(\mathrm{def})}$ is common to a specified stable defect class. Otherwise the mode-resolved notation $\rho_{\mathrm{IC}}^{(\mathrm{imp})}$ must be retained.

\section{Reduced Swept-Volume Approximation for Localized Defects}
\label{sec:AP_swept_volume}

The exact participation volume is
\begin{equation}
    V_{\mathrm{imp}} = \int_{\Omega_{\mathrm{IC}}^{(\mathrm{def})}} |\boldsymbol{\Upsilon}|^2\,\mathrm{d}V.
\end{equation}

For comparison with the thin-layer defect parametrization, consider a structural layer of equilibrium thickness $r_{\mathrm{amp}}$ surrounding a spherical Soft Floor node of radius $r_{\mathrm{D}}$. Its geometric shell volume is
\begin{equation}
    V_{\mathrm{sweep}} = \frac{4\pi}{3} \left[ (r_{\mathrm{D}}+r_{\mathrm{amp}})^3 - r_{\mathrm{D}}^3 \right].
\end{equation}

Expanding,
\begin{equation}
    V_{\mathrm{sweep}} = 4\pi r_{\mathrm{D}}^2 r_{\mathrm{amp}} \left[ 1 + \frac{r_{\mathrm{amp}}}{r_{\mathrm{D}}} + \frac{1}{3} \left( \frac{r_{\mathrm{amp}}}{r_{\mathrm{D}}} \right)^2 \right].
\end{equation}

Hence, for $r_{\mathrm{amp}} \ll r_{\mathrm{D}}$,
\begin{equation}
    V_{\mathrm{sweep}} = 4\pi r_{\mathrm{D}}^2 r_{\mathrm{amp}} \left[ 1 + \mathcal{O}\left(\frac{r_{\mathrm{amp}}}{r_{\mathrm{D}}}\right) \right].
\end{equation}

Define the leading geometric shell volume by
\begin{equation}
    V_{\mathrm{eff}} \equiv 4\pi r_{\mathrm{D}}^2 r_{\mathrm{amp}}.
\end{equation}

The kinetic participation volume need not equal this geometric volume. Define
\begin{equation}
    \mathcal{G}_{\mathrm{imp}} \equiv \frac{V_{\mathrm{imp}}}{V_{\mathrm{sweep}}}.
\end{equation}

The value of $\mathcal{G}_{\mathrm{imp}}$ is determined by the solved defect mode and may depend on geometry, boundary kinematics, and $r_{\mathrm{amp}}/r_{\mathrm{D}}$. A controlled thin-layer reduction requires a finite limiting value of $\mathcal{G}_{\mathrm{imp}}$ for the selected defect family.

Then
\begin{equation}
    V_{\mathrm{imp}} = 4\pi\mathcal{G}_{\mathrm{imp}} r_{\mathrm{D}}^2 r_{\mathrm{amp}} \left[ 1 + \mathcal{O}\left(\frac{r_{\mathrm{amp}}}{r_{\mathrm{D}}}\right) \right].
\end{equation}

The closure $\mathcal{G}_{\mathrm{imp}}=1$ is therefore mode-specific rather than geometric.

Using $m_{\mathrm{add}} = \rho_{\mathrm{IC}}^{(\mathrm{imp})}V_{\mathrm{imp}}$,
\begin{equation}
    m_{\mathrm{add}} = 4\pi r_{\mathrm{D}}^2 \rho_{\mathrm{IC}}^{(\mathrm{imp})} \mathcal{G}_{\mathrm{imp}} r_{\mathrm{amp}} \left[ 1 + \mathcal{O}\left(\frac{r_{\mathrm{amp}}}{r_{\mathrm{D}}}\right) \right].
\end{equation}

Define
\begin{equation}
    K_{\mathrm{D}}^{(\mathrm{def})} \equiv 4\pi r_{\mathrm{D}}^2 \rho_{\mathrm{IC}}^{(\mathrm{imp})} \mathcal{G}_{\mathrm{imp}} = 4\pi r_{\mathrm{D}}^2 \mu_{\mathrm{D}} \iota_{\mathrm{V}}^{(\mathrm{def})} \mathcal{G}_{\mathrm{imp}}.
\end{equation}

Then
\begin{equation}
    m_{\mathrm{add}} = K_{\mathrm{D}}^{(\mathrm{def})} r_{\mathrm{amp}} \left[ 1 + \mathcal{O}\left(\frac{r_{\mathrm{amp}}}{r_{\mathrm{D}}}\right) \right],
\end{equation}
with
\begin{equation}
    [K_{\mathrm{D}}^{(\mathrm{def})}] = \mathrm{kg\,m^{-1}}.
\end{equation}

The reduced relation $K_{\mathrm{D}} = 4\pi r_{\mathrm{D}}^2\rho_{\mathrm{IC}}$ is recovered only when the projected impedance density is universal within the selected defect class and $\mathcal{G}_{\mathrm{imp}}=1$.

\section{Phase-Participation Function and Branch-Level Sine Closure}
\label{sec:AP_phase_participation}

The localized boundary-value problem fixes the hydrodynamic inertia of a prescribed reference defect mode. A multiplicative phase law is admissible only when the phase dependence can be separated from the reference geometry and normalized flow field.

Under that factorization, introduce the nonnegative periodic participation function
\begin{equation}
    0 \leq f_{\varphi}(\Delta\varphi) \leq 1.
\end{equation}

For a spatially uniform phase mismatch,
\begin{equation}
    T_{\mathrm{def}} = \frac{1}{2} \rho_{\mathrm{IC}}^{(\mathrm{imp})} V_{\mathrm{imp}} f_{\varphi}(\Delta\varphi) \dot{q}^2,
\end{equation}
so that
\begin{equation}
    M_{\mathrm{top}}(\Delta\varphi) = \rho_{\mathrm{IC}}^{(\mathrm{imp})} V_{\mathrm{imp}} f_{\varphi}(\Delta\varphi).
\end{equation}

The factorization must not double-count phase dependence. If $\mathbf{b}_{\partial}$, $\Omega_{\mathrm{def}}$, or $\boldsymbol{\Upsilon}$ already depends explicitly on $\Delta\varphi$, the phase-dependent boundary-value functional is primary and an additional multiplicative $f_{\varphi}$ may be introduced only after a nonredundant separation is established.

The participation function satisfies
\begin{equation}
    f_{\varphi}(\Delta\varphi+2\pi) = f_{\varphi}(\Delta\varphi).
\end{equation}

For an unpolarized reciprocal defect sector,
\begin{equation}
    f_{\varphi}(-\Delta\varphi) = f_{\varphi}(\Delta\varphi),
\end{equation}
whereas an individual chiral branch need not be even.

For a phase-created defect in the factorized description,
\begin{equation}
    f_{\varphi}(0) = 0.
\end{equation}

If a geometric boundary persists at phase matching, its phase-independent added inertia must instead be retained separately,
\begin{equation}
    M_{\mathrm{add}}(\Delta\varphi) = M_{\mathrm{geom}} + M_{\varphi}f_{\varphi}(\Delta\varphi).
\end{equation}

The sine closure is admissible branchwise. On the positive branch,
\begin{equation}
    f_{\varphi}(\Delta\varphi) = \sin(\Delta\varphi), \qquad 0 \leq \Delta\varphi \leq \pi,
\end{equation}
while on the reciprocal negative branch,
\begin{equation}
    f_{\varphi}(\Delta\varphi) = -\sin(\Delta\varphi), \qquad -\pi \leq \Delta\varphi \leq 0.
\end{equation}

The reciprocal nonnegative envelope is therefore
\begin{equation}
    f_{\varphi}^{(\mathrm{rec})}(\Delta\varphi) = |\sin(\Delta\varphi)|.
\end{equation}

The sign of $\sin(\Delta\varphi)$ may encode orientation in a chiral force or circulation sector, but it cannot make the quadratic inertial coefficient negative.

For a nonuniform phase mismatch, the phase factor remains inside the weighted boundary integral,
\begin{equation}
    T_{\mathrm{def}} = \frac{1}{2}\mu_{\mathrm{D}}\dot{q}^2 \int_{\partial\Omega_{\mathrm{def}}} |\mathbf{b}_{\partial}\cdot\hat{\mathbf{n}}_{\mathrm{d}}|^2 f_{\varphi}[\Delta\varphi(\boldsymbol{s})] \,\mathrm{d}A.
\end{equation}

The scalar factorization therefore assumes a phase mismatch that is uniform, or effectively uniform under the normal-mode weighting, over the active interface.

\section{Recovery Conditions for the Hydrodynamic Mass Law}
\label{sec:AP_recovery_mass_law}

For a factorized reference defect mode,
\begin{equation}
    M_{\mathrm{top}} = \rho_{\mathrm{IC}}^{(\mathrm{imp})} V_{\mathrm{imp}} f_{\varphi}(\Delta\varphi).
\end{equation}

Using the thin-layer reduction,
\begin{equation}
    M_{\mathrm{top}} = K_{\mathrm{D}}^{(\mathrm{def})} r_{\mathrm{amp}} f_{\varphi}(\Delta\varphi) \left[ 1 + \mathcal{O}\left(\frac{r_{\mathrm{amp}}}{r_{\mathrm{D}}}\right) \right].
\end{equation}

The historical hydrodynamic mass law is recovered at leading order under
\begin{equation}
    r_{\mathrm{amp}} \ll r_{\mathrm{D}}, \qquad \mathcal{G}_{\mathrm{imp}} = 1, \qquad \rho_{\mathrm{IC}}^{(\mathrm{imp})} = \rho_{\mathrm{IC}},
\end{equation}
together with a uniform factorized phase mismatch and the positive-branch closure
\begin{equation}
    f_{\varphi}(\Delta\varphi) = \sin(\Delta\varphi), \qquad 0 \leq \Delta\varphi \leq \pi.
\end{equation}

Here $\rho_{\mathrm{IC}}^{(\mathrm{imp})} = \rho_{\mathrm{IC}}$ is a universality closure for the projected defect impedance density; it does not identify $\rho_{\mathrm{IC}}$ with an independently derived intrinsic equilibrium condensate density.

With
\begin{equation}
    V_{\mathrm{eff}} = 4\pi r_{\mathrm{D}}^2 r_{\mathrm{amp}},
\end{equation}
the leading reduction is
\begin{equation}
    M_{\mathrm{top}} = \rho_{\mathrm{IC}}V_{\mathrm{eff}} \sin(\Delta\varphi) \left[ 1 + \mathcal{O}\left(\frac{r_{\mathrm{amp}}}{r_{\mathrm{D}}}\right) \right].
\end{equation}

Defining
\begin{equation}
    K_{\mathrm{D}} = 4\pi r_{\mathrm{D}}^2\rho_{\mathrm{IC}},
\end{equation}
gives
\begin{equation}
    M_{\mathrm{top}} = K_{\mathrm{D}}r_{\mathrm{amp}} \sin(\Delta\varphi) \left[ 1 + \mathcal{O}\left(\frac{r_{\mathrm{amp}}}{r_{\mathrm{D}}}\right) \right].
\end{equation}

Thus the historical expression is a controlled thin-layer reduction of the mode-resolved boundary-value formulation rather than the unrestricted microscopic law.

For the reciprocal envelope,
\begin{equation}
    M_{\mathrm{top}} = \rho_{\mathrm{IC}}V_{\mathrm{eff}}|\sin(\Delta\varphi)|.
\end{equation}

When a persistent impermeable boundary remains at phase matching,
\begin{equation}
    M_{\mathrm{add}} = M_{\mathrm{geom}} + \rho_{\mathrm{IC}}^{(\mathrm{imp})} V_{\mathrm{imp}} f_{\varphi}(\Delta\varphi).
\end{equation}

This localized added inertia is distinct from the coherent primitive-cell scaffold inertia.

Under the TFT identification that the same mode-resolved topological loading serves as both the inertial coefficient and the coupling coefficient of the restorative tension potential,
\begin{equation}
    m_{\mathrm{i}} = M_{\mathrm{top}}, \qquad \mathbf{F}_{\mathrm{grav}} = -M_{\mathrm{top}}\nabla\Phi_{\tau},
\end{equation}
one obtains
\begin{equation}
    \mathbf{a} = \frac{\mathbf{F}_{\mathrm{grav}}}{m_{\mathrm{i}}} = -\nabla\Phi_{\tau}.
\end{equation}

This cancellation is conditional on using the same normalized loading functional in both the inertial and restorative couplings; it is not an independent derivation of that identification.

\section[Empirical Projection and Calibration]{Empirical Projection and Defect-Flow Calibration}
\label{sec:AP_empirical_projection}

The localized projection
\begin{equation}
    \iota_{\mathrm{V}}^{(\mathrm{def})} = \frac{\displaystyle\int_{\partial\Omega_{\mathrm{def}}} |\mathbf{b}_{\partial}\cdot\hat{\mathbf{n}}_{\mathrm{d}}|^2\,\mathrm{d}A} {\displaystyle\int_{\Omega_{\mathrm{IC}}^{(\mathrm{def})}} |\boldsymbol{\Upsilon}|^2\,\mathrm{d}V}
\end{equation}
is fully defined once the defect boundary, boundary mode, connected interstitial domain, and outer condition are specified. Its numerical value remains undetermined until the corresponding boundary-value problem is solved.

The assignment
\begin{equation}
    \iota_{\mathrm{V}} = 1\ \mathrm{m}^{-1}
\end{equation}
is therefore empirical. It corresponds to
\begin{equation}
    \ell_{\mathrm{imp}} = 1\ \mathrm{m},
\end{equation}
and gives
\begin{equation}
    \rho_{\mathrm{IC}}^{(\mathrm{emp})} = \mu_{\mathrm{D}} \left(1\ \mathrm{m}^{-1}\right).
\end{equation}

For $\mathcal{G}_{\mathrm{imp}}=1$,
\begin{equation}
    K_{\mathrm{D}}^{(\mathrm{emp})} = 4\pi r_{\mathrm{D}}^2 \mu_{\mathrm{D}} \left(1\ \mathrm{m}^{-1}\right).
\end{equation}

The particle-amplitude calibration depends explicitly on the adopted projection scale. With the existing interstitial-volume coefficient,
\begin{equation}
    r_{\mathrm{D}}(\iota_{\mathrm{V}}) = \left( \frac{Gm_{\mathrm{VEV}}} {5.8722\,c_{\mathrm{E}}^2\iota_{\mathrm{V}}} \right)^{1/2},
\end{equation}
so that
\begin{equation}
    r_{\mathrm{D}} \propto \iota_{\mathrm{V}}^{-1/2}.
\end{equation}

Consequently, calibrations using $\rho_{\mathrm{IC}}=\mu_{\mathrm{D}}\iota_{\mathrm{V}}$ inherit the adopted projection scale. A revised microscopic projection propagates into $r_{\mathrm{D}}$, $\mu_{\mathrm{D}}$, $\rho_{\mathrm{IC}}$, $K_{\mathrm{D}}$, $M_{\mathrm{top}}$, and quantities derived from them.

A complete microscopic calculation must solve the normal-flow problem in the connected FCC interstitial geometry with controlled outer conditions and demonstrate numerical convergence. Its conservative low-frequency limit must satisfy
\begin{equation}
    \lim_{\omega\rightarrow0} \iota_{\mathrm{V}}^{(\mathrm{def})}(\omega) = \iota_{\mathrm{V},0}^{(\mathrm{def})},
\end{equation}
provided the limit exists.

A universal projected impedance density may be defined only if this limit is common to a specified stable defect class,
\begin{equation}
    \iota_{\mathrm{V},0}^{(\mathrm{def})} = \iota_{\mathrm{V},0}.
\end{equation}

Then
\begin{equation}
    \rho_{\mathrm{IC}} \equiv \mu_{\mathrm{D}}\iota_{\mathrm{V},0}.
\end{equation}

Under this closure, $\rho_{\mathrm{IC}}$ is a universal projected impedance density for that defect class, not an independently established intrinsic equilibrium density of the condensate.

If a common geometry factor $\mathcal{G}_{0}$ also exists,
\begin{equation}
    K_{\mathrm{D}} = 4\pi r_{\mathrm{D}}^2 \mu_{\mathrm{D}} \iota_{\mathrm{V},0} \mathcal{G}_{0}.
\end{equation}

The historical relation $K_{\mathrm{D}}=4\pi r_{\mathrm{D}}^2\rho_{\mathrm{IC}}$ follows only for $\mathcal{G}_{0}=1$. Otherwise the defect-resolved quantities $\rho_{\mathrm{IC}}^{(\mathrm{imp})}$, $\iota_{\mathrm{V}}^{(\mathrm{def})}$, and $K_{\mathrm{D}}^{(\mathrm{def})}$ must be retained.

The final structure is therefore precise. The coherent branch satisfies $\mathbf{v}_{\mathrm{rel}}=\mathbf{0}$ and carries two finite-frequency transverse translational modes under the incremental volume constraint. Entrainment does not independently add condensate inertia to the primitive-cell scaffold normalization. The localized branch satisfies $\mathbf{v}_{\mathrm{rel}}\neq\mathbf{0}$ and acquires a conservative hydrodynamic added inertia determined by the defect boundary-value problem. The projection $\iota_{\mathrm{V}}^{(\mathrm{def})}$, its possible universality, and the corresponding canonical defect-flow solutions remain open microscopic calculations.
